\chapter{Reproducibility Details}
\label{app:method-details}

\section{Code and data files}
\label{app:schema}
\label{app:reproducibility}

The accompanying repository is at \url{https://github.com/D3prave/thesis}.

Most result tables and figures were generated from run-level measurements
for the short-answer experiment and the biography comparison in
Section~\ref{sec:longform-results}. Each measurement identifies one prompting
condition, dataset, generator, seed, temperature, entailment model,
correctness rule, and uncertainty method. Additional summaries describe the
variation across seeds and the agreement between the two LLM judges. The appendix tables
summarize these run values.

The hierarchical intervals of the short-answer comparisons come from a
record-level paired bootstrap analysis. Rebuilding tables from the resulting
summaries does not recompute the values from the raw answers.

Generation and LLM judging ran on the NHR@FAU Alex GPU cluster. Cross-encoder
clustering ran on Alex and on the Fritz CPU cluster. Given stored answers,
cross-encoder clustering and exact matching are deterministic. LLM entailment
decisions can vary with inference batching.

\section{Long-form biography records}
\label{app:longform-details}

The biography evaluation contains three generators, three seeds, and
500 subjects per generator--seed run. Prompts follow the FActScore subject
format~\cite{min_factscore_2023}. The system instruction requests a factual
biography of about 120--150 words, and the query is ``Tell me a bio of
\{entity\}.'' Ten responses are sampled using vLLM at $T=1.0$, top-$p=0.9$,
top-$k=50$, with at most 256 new tokens each. The saved Wikipedia introduction
provides the reference.

One paragraph per prompt is graded: the most frequent exact normalized
surface form, with ties resolved by first occurrence. Qwen2.5-72B-Instruct
returns one \texttt{CORRECT}/\texttt{INCORRECT} verdict for that paragraph.
Specific facts must be supported by the reference, while vague text without a
contradiction is permitted. The reported acceptance rate is the fraction of
selected paragraphs receiving \texttt{CORRECT}.

The recorded comparisons cluster normalized paragraphs into connected
components. DeBERTa uses answer-only input and requires entailment in both
directions. Qwen receives the question and asks whether both paragraphs convey
the same factual claims, requiring agreement in both response orders. This is
the biography pipeline's grouping rule.
The short-answer pipeline uses the raw-answer anchor scan in
Section~\ref{sec:clustering}. The four settings are exact matching,
DeBERTa-v3-base~\cite{sentence_transformers_nli_deberta_v3_base_2022},
DeBERTa-v3-large, and Qwen-72B. Sequence log-probabilities were not retained
because the vLLM path exposed tempered, truncated probabilities, so this
comparison uses count-based entropy only. Figure~\ref{fig:longform-comparison}
and Tables~\ref{tab:longform}--\ref{tab:longform-by-generator} report the results.

The records-only intervals were computed from the saved labeled biography
records using the same AUROC and paired bootstrap functions as the
short-answer analysis.

For each of 2{,}000 resamples, with random seed 20260831, the analysis draws
500 biographies with replacement within each generator--seed run. Both
estimators and the correctness label use the same sampled indices. Draws are
independent across runs. It averages the paired AUROC differences over the
three seeds, then equally over the three generators for the pooled result.
The 95\,\% intervals are the 2.5th and 97.5th percentiles of these differences.
All 108 relevant run-level AUROCs match the recorded measurements to six decimals.
No resample had an undefined AUROC.

The dataset, generators, and seeds remain fixed. With one dataset and three
dataset--generator units, no unit or hierarchical intervals were computed.
These are the permissive records-only intervals described in
Section~\ref{sec:statistics}, not estimates of uncertainty across
new datasets and generators. The evaluation also uses one correctness judge,
Qwen, which is one of the entailment models.

\section{Collection checks}
\label{app:collection-checks}

\paragraph{Temperature-sweep labels.}
LLM correctness labels are matched by prompt identifier to the main
collection. They are bit-identical across all 108 checked cells, so all sweep
temperatures are evaluated against the same temperature-$0.1$ answer.
Token-F1 instead grades the modal sampled answer because the sweep did not
save a separate most-likely answer. Its accuracy differs from the main run in
all 54 checked cells, by 0.015 on average in absolute terms and at most 0.262.
The reported sweep therefore uses only LLM correctness labels, and its
answer-length fields are not compared across temperatures.

\paragraph{TriviaQA deduplication.}
\label{app:dedup}
The SQuAD-format conversion stores multiple question--context rows for some
questions, and its row-level split places copies of 16 questions in both
training and evaluation. A further nine questions repeat within training and
six within evaluation. Removing duplicate questions leaves 391 training and
378 evaluation questions. The same evaluation exclusions apply to every
method, entailment model, grader, and record-level resampling analysis.
NQ-Open and SVAMP contain no repeats in the selected sets.

\section{Bootstrap details}
\label{sec:bootstrap-assumptions}

The saved short-answer paired analysis uses 2,000 resamples. Within each selected
dataset--generator pair, it averages the three seed-level AUROC differences,
then averages over the nine selected pairs. Record resampling keeps the
compared methods paired within a seed cell, but draws question indices
separately across seeds and generators. It neither synchronizes repeated
questions across these cells nor resamples datasets and generators as
separate factors.

A record resample with only one correctness class has undefined AUROC and is
omitted from that draw's average. This can matter when few answers are labeled
correct. Omission rates were not retained in the original export. The additional
analyses in Appendix~\ref{app:robustness} report them explicitly. The original
intervals and bootstrap tail
proportions are conditional on these choices and are not adjusted for
multiple comparisons. The released code instead reports 90\,\% bias-corrected
and accelerated intervals from 1{,}000 resamples, so its intervals are not
directly comparable with these.

Run-level measurements can verify paired point estimates, but do not suffice
to recompute the record-level intervals, which require the scored answers.
Appendix~\ref{app:robustness} reports additional record-level
analyses.

\section{Probe training and saved features}
\label{app:probe-training}

For each training question, the most-likely answer supplies a hidden state.
Ten separate sampled responses supply the discrete semantic entropy target.
The high/low threshold minimizes the sum of squared deviations within the
two entropy groups, following Equation~(5) of
Kossen et al.~\cite{kossen_semantic_entropy_probes_2024}. A logistic-regression
classifier predicts that label from the hidden state. The accuracy probe
instead learns correctness labels. At evaluation time, each probe needs one
answer and its hidden state.

Training is separate for each dataset, generator, and seed, using questions
disjoint from evaluation. Feature choices are fixed rather than tuned on the
evaluation set. The hidden state comes from position
$n_{\mathrm{generated}}-1$, where the stop sequence determines
$n_{\mathrm{generated}}$. An end-of-sequence token need not be present.

All 108 probes use final-layer features, last-token pooling,
$L_2$ regularization with $C=1$, a training seed of zero and at most 2{,}000
optimizer iterations, fitted with the \texttt{lbfgs} solver without feature
standardization. All fits report convergence. Feature dimensions are 4{,}096
for the two smaller models and 8{,}192 for Llama-3.1-70B.

The accuracy probe learns the stored training correctness labels. Scoring it
with another judge does not retrain it. The stored training records do not
say which judge produced these labels.

For the NQ-Open Llama-3.1-70B short-phrase run with seed 42, both probes report
399 training examples and one excluded prompt, which matches the one training
record without a saved feature. All evaluation records have saved features.
